%% ---------------------------------------------------------------------
\section{El decibel}
\label{sec:rf-decibel}
%% ---------------------------------------------------------------------

\subsection{Un invento de ingenieros telefónicos}

Antes de 1923, la pérdida de una línea telefónica se expresaba en \emph{millas de cable estándar} (MSC): cuántas millas de un cable de referencia atenuaban lo mismo a $5000$\,rad/s ($\approx796$\,Hz, una frecuencia representativa de la voz). La unidad dependía de la frecuencia y del cable. En 1924, W.~H.~Martin, de Bell Labs, propuso reemplazarla por la \emph{transmission unit} (TU): diez veces el logaritmo decimal de la razón de potencias~\citep{martin24}. La TU casi coincidía con la milla de cable,
\begin{equation}
  1\ \text{MSC}\approx\frac{1}{1.056}\ \text{TU}\approx 0.95\ \text{dB},
\end{equation}
de modo que los ingenieros no tuvieron que cambiar su intuición. En 1928 el Bell System adoptó el nombre \emph{decibel}, en honor a Alexander Graham Bell, siguiendo la propuesta internacional de llamar \emph{bel} a la unidad decimal y \emph{neper} (por John Napier) a la natural~\citep{martin29}:
\begin{gather}
  \text{bel}=\log_{10}\frac{P_2}{P_1},\qquad
  \text{dB}=10\log_{10}\frac{P_2}{P_1},\qquad
  1\ \text{Np}=20\log_{10}e=8.686\ \text{dB}.
\end{gather}
El bel resultó demasiado grande para el trabajo diario; sobrevivió su décima parte.

\subsection{¿Por qué 10 log y por qué 20 log?}

El decibel se define para \emph{potencias}. Si la potencia se disipa sobre la misma resistencia, $P=V^2/R$ y
\begin{equation}
  10\log_{10}\frac{V_2^2/R}{V_1^2/R}=20\log_{10}\frac{V_2}{V_1}.
\end{equation}
En el diagrama de Bode $|\Gma(\jw)|$ es una razón de \emph{amplitudes} (salida sobre entrada); por eso siempre se usa
\begin{equation}
  |\Gma|_{\mathrm{dB}}=20\log_{10}|\Gma(\jw)|.
\end{equation}

\begin{table}[ht]
\centering\small
\begin{tabular}{@{}cc l@{}}
\toprule
$|G|$ & dB & para recordar\\
\midrule
$10$ & $+20$ & una década de amplitud\\
$2$ & $+6.02$ & el doble\\
$\sqrt2$ & $+3.01$ & la esquina de un cero real\\
$1$ & $0$ & cruce de ganancia\\
$1/\sqrt2$ & $-3.01$ & ancho de banda; esquina de un polo real\\
$0.1$ & $-20$ & una década menos\\
\bottomrule
\end{tabular}
\caption{Valores que conviene memorizar. La mitad de la \emph{potencia} es $-3$\,dB: por eso la frecuencia de corte de un filtro se define donde la magnitud cae $3$\,dB.}
\label{tab:rf-db}
\end{table}

\paragraph{Tres razones para usar logaritmos}
\begin{enumerate}
  \item \textbf{Cascadas.} Las ganancias se multiplican; en dB se suman, y las fases también se suman:
  \[
    20\log|G_1G_2G_3|=20\log|G_1|+20\log|G_2|+20\log|G_3|,\qquad
    \angle(G_1G_2G_3)=\angle G_1+\angle G_2+\angle G_3 .
  \]
  Por ejemplo, $\times10\cdot\times0.5\cdot\times4=\times20$ en amplitud equivale a $20-6.02+12.04=26.02$\,dB.
  \item \textbf{Rangos enormes.} Una razón de amplitudes de $10^6$ son ``solo'' $120$\,dB.
  \item \textbf{Percepción.} La sonoridad que percibimos crece aproximadamente con el logaritmo de la intensidad; un cambio de $1$\,dB apenas se nota.
\end{enumerate}

\notamargen{Cuando se agrega una referencia, el decibel deja de ser una razón y pasa a ser un \emph{nivel}: dBm ($1$\,mW), dBV ($1$\,V), dB SPL ($20\,\mu$Pa). Un número en dB sin referencia es siempre una razón.}

\begin{pausa}
Un amplificador de $+26$\,dB se conecta en cascada con un filtro que, a cierta frecuencia, atenúa $-32$\,dB. ¿Por cuánto queda multiplicada la amplitud de una senoide de esa frecuencia?
\respuesta{$26-32=-6$\,dB, es decir, $10^{-6/20}=0.50$: la amplitud queda a la mitad. En dB las cascadas se suman.}
\end{pausa}
